% Slide 5: Spatial Discretization: CAD-IGA Integration
\begin{frame}{Spatial Discretization: CAD-IGA Integration}
    \begin{itemize}
        \item \textbf{Isogeometric Analysis (IGA)}: Uses Non-Uniform Rational B-Splines (NURBS) from CAD directly for physics:
        \begin{equation*}
            R_{i,j}(\xi, \eta) = \frac{N_{i,p}(\xi) M_{j,q}(\eta) w_{i,j}}{\sum_{k,l} N_{k,p}(\xi) M_{l,q}(\eta) w_{k,l}}
        \end{equation*}
        \item \textbf{Exact Geometry}: Smooth notches and circular boundaries are represented exactly at any discretization level.
        \item \textbf{Higher Continuity}: $C^{p-1}$ continuity across knots reduces shear locking, ensuring superior accuracy per degree of freedom compared to Lagrangian FEM.
    \end{itemize}
    \centering
    \vspace{0.1cm}
    \begin{tikzpicture}[scale=0.7]
        % Show control net mapping
        \draw[gray, dotted] (0,0) grid (4,1);
        \draw[draw=navyblue, thick] (0,0.5) to[out=30,in=150] (4,0.5);
        \node[font=\tiny, left] at (0,0.5) {NURBS Control Net};
        \fill[red] (1,0.65) circle (2pt) node[above, font=\tiny] {$P_{i,j}$};
    \end{tikzpicture}
\end{frame}

% Slide 6: Mesh Convergence Study
\begin{frame}{Mesh Convergence \& Justification}
    To justify the $80 \times 40$ baseline mesh resolution, a static convergence study was performed:
    \begin{table}
        \centering
        \resizebox{0.8\textwidth}{!}{%
        \begin{tabular}{lcccc}
            \toprule
            \textbf{Grid Resolution} & \textbf{Elements} & \textbf{DoFs} & \textbf{Tip Displacement $U_y$ (mm)} & \textbf{Relative Error (\%)} \\
            \midrule
            $10 \times 5$   & 50   & 168   & 3.538 & 16.51\% \\
            $20 \times 10$  & 200  & 528   & 3.889 & 8.24\% \\
            $40 \times 20$  & 800  & 1,848 & 4.083 & 3.66\% \\
            $\mathbf{80 \times 40}$ & \textbf{3,200} & \textbf{6,888} & \textbf{4.185} & \textbf{1.24\% (Selected)} \\
            $160 \times 80$ & 12,800 & 26,568 & 4.238 & Reference \\
            \bottomrule
        \end{tabular}%
        }
    \end{table}
    \begin{itemize}
        \item Refining to $160 \times 80$ quadruples the element count and degree of freedom count ($26,568$), but only shifts the displacement by $1.24\%$.
        \item The $80 \times 40$ mesh represents the optimal trade-off point between numerical convergence and CPU assembly cost.
    \end{itemize}
\end{frame}
